% Penutup buku. Bagian ini memuat kesimpulan yang dituntut sebuah buku
% referensi, yaitu sintesis, jawaban atas rumusan masalah pada Pendahuluan,
% kontribusi, keterbatasan, dan arah pengembangan. Seluruh angkanya dikutip
% dari Summary dan tabel contoh bab masing-masing.
\chapter{Penutup}
\label{ch:penutup}

% Penutup tidak bernomor bab. Pencacah bab dinolkan agar nomor tabelnya tidak
% berawalan 15, lalu nomornya melanjutkan tabel terakhir Pendahuluan.
\setcounter{chapter}{0}
\setcounter{table}{\value{tabelluarbab}}

\section*{Sintesis Isi Buku}
\addcontentsline{toc}{section}{Sintesis Isi Buku}

Keempat belas bab buku ini membahas gaya yang berbeda, tetapi polanya sama.
Setiap gaya memperoleh satu sifat dengan membayar sifat lain, dan pembayaran itu
dapat diukur. \textit{Layered architecture} menahan perubahan di satu
\textit{layer} dengan membayar \textit{hop} tambahan. \textit{Event-driven
architecture} melepas \textit{publisher} dari nasib \textit{subscriber} dengan
membayar latensi. Command Query Responsibility Segregation (CQRS) memisahkan
model tulis dari model baca dengan membayar dua kali simpan. \textit{Serverless}
menekan biaya saat diam dengan membayar \textit{cold start}.

Tabel~\ref{tab:harga-gaya} merangkum pasangan tersebut untuk setiap bab. Kolom
ketiga memuat yang diperoleh, sedangkan kolom keempat memuat harga yang dibayar.
Keduanya dinyatakan dengan angka dari pengukuran pada bab masing-masing,
sehingga setiap baris dapat diperiksa ulang lewat Latihan bab tersebut.

\begin{table}[htbp]
  \centering
  \caption{Yang diperoleh dan harga yang dibayar setiap gaya, menurut
    pengukuran pada babnya}
  \label{tab:harga-gaya}
  \begin{tabularx}{\textwidth}{|c|>{\raggedright\arraybackslash}p{1.9cm}%
    |>{\raggedright\arraybackslash}X|>{\raggedright\arraybackslash}X|}
    \hline
    \textbf{Bab} & \textbf{Gaya} & \textbf{Yang diperoleh} &
      \textbf{Harga yang dibayar} \\
    \hline
    1 & Client-server & Data bersama dan aturan bisnis berada di satu
      server & Antrean: median hanya naik ke 1,36 ms pada enam belas klien,
      tetapi nilai terbesarnya melampaui 1.000 ms sejak delapan klien \\
    \hline
    2 & \textit{Layered architecture} & Perubahan satu urusan teknis tertahan
      di satu \textit{layer} & Satu \textit{hop} tambahan bernilai 0,0046 ms
      pada median, sekitar 5,5 persen satu \textit{request} \\
    \hline
    3 & Model-View-* & Model-View-Intent (MVI) mencatat nol \textit{dependency}
      dan diuji tanpa View & Setiap tindakan pengguna harus dibungkus menjadi
      Intent, sehingga kodenya bertambah \\
    \hline
    4 & \textit{Ports and adapters} & Berpindah dari memori ke SQLite mengubah
      nol \textit{file} pada \textit{core logic} & Setiap kebutuhan keluar
      menambah \textit{port} dan \textit{adapter}; \textit{adapter} ketiga
      menambah dua \textit{file} dan 29 baris \\
    \hline
    5 & Microkernel & Nol penyebutan nama \textit{plugin} di \textit{core} &
      \textit{Contract}, \textit{manifest}, dan \textit{registry}; tanpa
      \textit{error handler}, satu \textit{plugin} rusak menghentikan
      aplikasi \\
    \hline
    6 & \textit{Pipe-and-filter} & Nol penyebutan antar \textit{filter} dari
      lima \textit{filter} & Penukaran \textit{stage} yang sah belum tentu
      menghasilkan \textit{file} yang sama \\
    \hline
    7 & \textit{Peer-to-peer} & Giliran terburuk saat satu \textit{node} mati
      55,2 dan 71,5 persen, melawan nol persen pada client-server & 19 dan
      empat \textit{message} per \textit{lookup}, melawan dua pada
      \textit{centralized index} \\
    \hline
    8 & \textit{Space-based architecture} & 23.450,0 \textit{request} per
      detik pada satu unit, melawan 9.194,9 pada versi basis data & Selisih
      48 kuota pada lag 100 ms, 222,0 tabrakan per detik, serta 8 sampai 45
      \textit{request} hilang setiap kali satu unit mati pada penyangga lima
      puluh \\
    \hline
    9 & \textit{Event-driven architecture} & \textit{Publisher} kembali sesudah
      0,1 ms, melawan \textit{blocking} 513,1 ms untuk seratus
      \textit{request} pada susunan sinkron & Latensi ujung ke ujung 6,2 kali
      pada topologi broker dan 9,3 kali pada topologi mediator \\
    \hline
    10 & CQRS dan \textit{event sourcing} & Pertanyaan riwayat yang terjawab
      naik dari nol menjadi satu pada CQRS, lalu dua pada \textit{event
      sourcing} & Latensi tulis CQRS 56,2983 ms melawan 28,0638 ms pada versi
      yang menimpa baris;
      \textit{file event sourcing} 55,7 kali \textit{file} keadaan \\
    \hline
    11 & \textit{Service-based architecture} & Saat bidang kurs mati, dua dari
      empat \textit{request} tetap terlayani, melawan nol pada monolit &
      Median bertambah 4,68 ms, sekitar 1,56 ms untuk setiap \textit{hop}
      tambahan \\
    \hline
    12 & Orkestrasi & Menyisipkan satu langkah cukup menyunting satu
      \textit{file}, melawan urutan yang tersebar di enam \textit{file} &
      Selisih median 0,478 ms; perubahan satu layanan perusahaan menuntut
      empat dari lima alur kerja diuji ulang \\
    \hline
    13 & Mikroservis & Rilis ulang kurs tidak menggagalkan satu pun
      \textit{request}, melawan 7 dari 200 pada monolit & Laporan 4,8 kali
      lebih lambat, dengan \textit{request} jaringan naik dari satu menjadi
      lima \\
    \hline
    14 & \textit{Serverless} & Saat diam memegang 16,0 KiB, melawan 21.179,3
      KiB pada layanan yang hidup terus & \textit{Cold start} 120,5837 ms
      melawan \textit{warm start} 0,0059 ms; titik impas 62.069
      \textit{request} per jam \\
    \hline
  \end{tabularx}
\end{table}

Pelajaran berikutnya menyangkut cara membaca angka. Rata-rata sering
menyembunyikan yang penting. Pada Bab~1, median waktu tanggap hanya naik dari
0,31 menjadi 1,36 milidetik, sedangkan nilai terbesarnya melampaui seribu
milidetik. Pada Bab~7, rata-rata \textit{lookup} yang terjawab saat satu
\textit{node} mati hampir sama pada ketiga susunan, yaitu 83 sampai 88 persen.
Giliran terburuknya justru berkisar dari nol sampai 71,5 persen.

Harga terbesar juga jarang muncul pada keadaan biasa. Harga itu muncul ketika
sebuah \textit{node} mati pada Bab~7, ketika \textit{validator} tersendat pada
Bab~9, dan ketika \textit{request} datang serentak sesudah \textit{scale to zero}
pada Bab~14. Pengukuran yang hanya menjalankan keadaan biasa karena itu tidak
akan menemukan harga tersebut.

\section*{Jawaban atas Rumusan Masalah}
\addcontentsline{toc}{section}{Jawaban atas Rumusan Masalah}

Pertanyaan pertama menanyakan cara merumuskan klaim kualitas menjadi
\textit{falsifiable claim}. Jawabannya, klaim dirumuskan sebagai hitungan yang
dapat diambil dari kode. Klaim \textit{core independence} pada Bab~4 menjadi
jumlah teknologi \textit{storage} yang disebut \textit{core logic}, sedangkan
pada Bab~5 menjadi jumlah nama \textit{plugin} yang disebut \textit{core}.
Keduanya bernilai nol. Klaim \textit{filter independence} pada Bab~6 menjadi
jumlah penyebutan antar \textit{filter}, dan angkanya nol dari lima
\textit{filter}.

Klaim yang menyangkut waktu dirumuskan sebagai selisih yang menyebut arah dan
besarnya, lalu diperiksa dengan pengulangan. Bab~9 memakai bentuk itu untuk
anggapan bahwa event mempercepat sebuah \textit{request}. Anggapan itu gugur,
sebab susunan sinkron menang pada 200 dari 200 putaran.

Pertanyaan kedua menanyakan besarnya harga. Tabel~\ref{tab:harga-gaya} memuat
jawabannya untuk setiap bab. Sebagian besar harga itu dinyatakan dalam satuan
yang dapat dihitung, yaitu milidetik, jumlah \textit{message}, jumlah
\textit{file}, atau ukuran memori. Satuan seperti itulah yang membuat
\textit{trade-off} dapat ditimbang, bukan sekadar disebut.

Pertanyaan ketiga menanyakan kapan selisih cukup kuat untuk menentukan pilihan.
Buku ini memakai dua uji berurutan. Uji pertama membandingkan selisih median
dengan \textit{range}. Pada Bab~5, selisih \textit{startup time} lebih besar
daripada \textit{range}-nya, sehingga selisih itu langsung dapat disebut ada.

Uji kedua dipakai ketika selisihnya lebih kecil daripada \textit{range}, yaitu
\textit{win count} pada seratus putaran atau lebih. Bab~2 mencatat 97 dari 100
putaran, sedangkan Bab~11 mencatat 140 dari 150 putaran. Pada kedua bab itu
selisih median lebih kecil daripada \textit{range}, tetapi arahnya tetap pasti.

Selisih yang sudah pasti ada belum tentu menentukan pilihan. Pada Bab~2,
penghematan \textit{open layer} terbukti pada 97 dari 100 putaran, tetapi dua
pemeriksaan yang dilepasnya hilang pada setiap \textit{request}. Keputusannya
karena itu kembali pada tuntutan bisnis setiap \textit{request}. Jalan pintas
dibenarkan untuk kurs dasar yang diterbitkan untuk umum, tetapi tidak untuk kurs
jual yang hanya boleh dibaca staf.

\section*{Kontribusi}
\addcontentsline{toc}{section}{Kontribusi}

Sumbangan utama buku ini adalah aturan pengukuran yang sama untuk empat belas
gaya. Aturan itu mewajibkan pengulangan sedikitnya seratus putaran untuk
selisih yang kecil, pelaporan median, persentil ke-95, dan nilai terbesar, serta
perbandingan selisih dengan \textit{range}. Hasil antarbab karena itu dapat
dibandingkan polanya.

Setiap angka juga disertai cara mengulanginya, yaitu program di
\berkas{sources/} dan 42 Latihan beserta tabel contohnya. Istilahnya dibakukan
dan dipakai seragam di seluruh bab, lalu dikumpulkan pada Glosarium.

\section*{Keterbatasan}
\addcontentsline{toc}{section}{Keterbatasan}

Empat keterbatasan perlu diingat ketika angka buku ini dipakai.

Pertama, seluruh pengukuran berjalan pada satu laptop dan alamat lokal. Waktu
tempuh jaringan antarkota tidak termuat, sehingga harga \textit{hop} pada Bab~11
sampai Bab~13 cenderung lebih kecil daripada pada jaringan sungguhan.

Kedua, skala contohnya kecil. \textit{Processing unit} pada Bab~8 diukur pada
satu, empat, dan tujuh unit di satu komputer. Angkanya memperlihatkan arah
pertumbuhan, sedangkan kapasitas sistem produksi perlu diukur pada sistem itu
sendiri.

Ketiga, angka waktu bergantung pada Python dan penerjemahnya. Pada Bab~9,
perpindahan antar \textit{thread} jauh lebih mahal daripada pekerjaan kecil yang
diangkutnya. Perbandingan itu dapat berbeda pada bahasa pemrograman lain.

Keempat, sebagian keadaan ditiru, bukan diambil dari layanan sungguhan. Penjalan
fungsi pada Bab~14 meniru tiga fase AWS Lambda di komputer lokal. Angka
\textit{cold start} dan gigabita detiknya karena itu mengikuti tiruan tersebut,
sehingga tagihan penyedia sungguhan dapat berbeda.

\section*{Arah Pengembangan}
\addcontentsline{toc}{section}{Arah Pengembangan}

Arah pertama adalah mengulang pengukuran yang sama pada jaringan sungguhan,
yaitu layanan yang berjalan di beberapa komputer. Langkah itu menguji apakah
pola pada Bab~11 sampai Bab~13 bertahan ketika harga satu \textit{hop} naik.

Arah kedua sudah dirancang di repositori buku ini, pada direktori
\berkas{papers/}. Sepuluh rancangan penelitian di sana melanjutkan bab-bab buku
ini, dan tujuh di antaranya menyangkut sistem yang memanggil model bahasa. Dua
rancangan mengukur
\textit{ports and adapters} ketika model bahasa berada di balik \textit{port},
termasuk biaya mengganti penyedianya. Dua rancangan menguji \textit{layered
architecture} ketika pemanggilan model dikumpulkan di satu \textit{layer} atau
dipanggil langsung dari tampilan.

Dua rancangan lain menerapkan topologi broker dan topologi mediator pada sistem
multi-agen berbasis model bahasa. Dari empat rancangan sisanya, dua menyangkut
gaya itu sendiri, yaitu pengenalan gaya arsitektur dari kode sumber dan pemetaan
gaya pada sistem pemerintahan digital di Indonesia. Dua lainnya menggabungkan
gaya, yaitu \textit{filter} yang dapat ditukar saat program berjalan dan
\textit{port} pada batas tampilan.

Pesan utama buku ini sederhana. Gaya arsitektur dipilih dengan membaca harganya,
dan harga itu baru dapat dibaca sesudah diukur.
